\documentclass[a4paper,14pt]{article}
\usepackage[utf8]{inputenc}
\usepackage[russian]{babel} 
\usepackage{graphicx} 
\usepackage{amsmath}
\usepackage{amssymb}
\usepackage{siunitx}% Для погрешностей
\usepackage{float}
\usepackage[colorlinks=true, linkcolor=blue]{hyperref}
\usepackage{enumitem}
\usepackage{longtable}
%\documentclass[fleqn]{article}

\renewcommand{\floatpagefraction}{0.99}
\renewcommand{\topfraction}{0.99}

% Подключаем пакет для управления колонтитулами
\usepackage{scrlayer-scrpage}
\pagestyle{scrheadings}

%\clearpagemainheadfoot % Очищаем стандартные колонтитулы
\cfoot{\thepage}
% Настраиваем верхний правый угол (Right Head)
\rofoot{\includegraphics[width=2.5cm]{Mir_insignia.png}}
% Вместо \rohead можно использовать:
% \lehead (левый верхний), \rofoot (правый нижний), \lefoot (левый нижний)

\title{Проект на тему "Европоцентризм"}


\author{Батанин Артём Александрович}
\date{\today}


\begin{document}

\maketitle

\newpage
\tableofcontents
\newpage



\section{Цель работы}


\section{История возникновения}

\subsection{Предпосылки}
    Европпейские страные к концу средних веков стали лидерами научного развития. Во время ВГО Европейская цивилизация
была мировым гегемоном в экономике, культуре и политике. Европоцентризм возник как ответ на вопрос - почему Европе удалось достичь
превосходства. Основная идея: западный образ жизни и культура сами способствуют прогрессу, в то время как остальные культуры являются архаизмами,
    тормозящими развитие.


\section{В России}
Россию нельзя однозначно отнести к Европе или к Востоку - она находится на границе двух миров. Это проявляется и
в разделении мнений граждан России. Россия до 1917 года - была часью западной цивилизации, (она была частью Антанты),
но после революции Советская Россия выбрала свой новый, особенный путь развития.

\section{Сторонники}


\section{Противники}

\section{Контрпримеры}

153

15
\end{document}